\section{Landau 理论的微观导出}

Gross 第 33 章从两点顶点的 $\omega$/$q$ 极限导出 Landau $f$ 函数，使唯象论与微观 $G$、$\Gamma$ 接通。

\subsection{前向散射}

准粒子之间的有效相互作用来自四点顶点在费米面、动量转移 $q\to0$ 的极限。两个极限不可交换：
\begin{equation}
  \Gamma^\omega
  =\lim_{\omega\to0}\lim_{q\to0}\Gamma_{\mathrm{ph}}(q,\omega),
  \qquad
  \Gamma^q
  =\lim_{q\to0}\lim_{\omega\to0}\Gamma_{\mathrm{ph}}.
\end{equation}
$\Gamma^\omega$ 对应准粒子尚未重新分布（Landau $f$ 的定义）；$\Gamma^q$ 对应热力学屏蔽后的响应。BS 方程给出
\begin{equation}
  \Gamma^q
  =\frac{\Gamma^\omega}{1+N(0)\langle\Gamma^\omega\rangle_{\mathrm{ang}}},
\end{equation}
（示意，$\ell=0$ 通道。）于是
\begin{equation}
  f_{\mathbf{k}\sigma,\mathbf{k}'\sigma'}
  =Z_{\mathbf{k}}Z_{\mathbf{k}'}
  \Gamma^\omega(\mathbf{k}\sigma,\mathbf{k}'\sigma'),
\end{equation}
$Z$ 来自单粒子极点留数。勒让德展开后即 $F_\ell^{s,a}$。

\subsection{有效质量}

微观
\begin{equation}
  \frac{m^*}{m}
  =\frac{1-\partial_\omega\mathrm{Re}\Sigma}{1+\partial_k\mathrm{Re}\Sigma/v_F^{(0)}}
  \Big|_{k_F}.
\end{equation}
伽利略不变把 $\partial_k\Sigma$ 与 $F_1^s$ 锁定，回到 $m^*/m=1+F_1^s/3$。这是 Ward 恒等式在 boost 下的结果。

\subsection{从微观到输运}

Gross 第 32 章：准粒子玻尔兹曼方程
\begin{equation}
  \partial_t\delta n_{\mathbf{k}}
  +\mathbf{v}_{\mathbf{k}}^*\cdot\nabla_{\mathbf{r}}
  \bigl(\delta n_{\mathbf{k}}-\frac{\partial n^0}{\partial\varepsilon}\delta\varepsilon_{\mathbf{k}}\bigr)
  =I_{\mathrm{coll}}[\delta n],
\end{equation}
$\delta\varepsilon=\sum f\delta n$。无碰撞极限给出零声；碰撞积分给出第一声衰减与粘滞。集体模式与第 5 章 $\chi(q,\omega)$ 的极点一致。
